\documentclass[10pt,a4paper]{amsart}
\usepackage[margin=19mm]{geometry}
\usepackage{amsmath,amssymb,mathtools}
\usepackage[scr=rsfs,cal=euler]{mathalfa}
\usepackage{xfrac}
\usepackage[unicode,hidelinks,bookmarks=false]{hyperref}
\newcommand{\ie}{i.e.\ }
\newcommand{\cf}{cf.\ }
\newcommand{\resp}{respectively\ }
\newcommand{\Subsection}{\subsection}
\providecommand{\nobreakdash}{\nobreak\hbox{-}}
\setlength{\emergencystretch}{2em}
\multlinegap0pt
\usepackage{amssymb}
\usepackage{tikz}
\usepackage{bm}
\newtheorem{theorem}{Theorem}
\newtheorem{lemma}{Lemma}
\newcommand{\Syl}{\mathop{\mathrm{Syl}}}
\title{Sylow subgroups for distinct primes\\and intersection of nilpotent subgroups}
\author{Francesca Lisi, Luca Sabatini}
\date{Source: arXiv:2505.21222v3 (CC BY 4.0)}
\begin{document}
\maketitle

\noindent\emph{Source and licence.} Sylow subgroups for distinct primes\\and intersection of nilpotent subgroups. Version arXiv:2505.21222v3 (CC BY 4.0), distributed by arXiv under the Creative Commons Attribution 4.0 International licence, http://creativecommons.org/licenses/by/4.0/, as recorded in the arXiv licence metadata for this exact version and shown on the abstract page https://arxiv.org/abs/2505.21222v3. Full licence notice: \texttt{licenses/arxiv-2505.21222-CC-BY-4.0.txt}.\par\smallskip

\noindent\textit{Editorial scope.} Counted unit: the article's lemma lem:key (source0.tex lines 661--666). The article's complete original proof of it is delivered with it (source0.tex lines 667--693, one \texttt{proof} environment of 27 lines). No mathematics is written, repaired or re-derived: the statement and the proof are the article's own lines, in the article's own notation, and no retained line is substituted or re-flowed.  Retained uncounted as the source-local context the proof needs: lines 378--381.  Nothing was omitted from any retained span, so the package carries no omission record.  No retained line contains a citation, so the wrapper carries no bibliography and no bibliography entry.  No retained line contains a figure environment, an image-inclusion command or drawing code, so no image is delivered and the package declares no asset dependency.  Editorial formatting only, in addition: an AMS \texttt{amsart} preamble that restates the article's own declarations and macros used by the retained lines, namely the environments \texttt{theorem}, \texttt{lemma} on separate counters, together with the standard packages those lines need.  The retained environments print in this document as Theorem 1, Lemma 1 in order of appearance, whereas the article prints the same items with the numbering of its own full document.  The complete native source of the article as distributed in the arXiv e-print for this exact version is bundled unchanged as \texttt{sources/178851/source0.tex}.
\begin{theorem}[Bialostocki] \label{thBialoHall}
Let $G$ be a group of odd order.
If $H \leqslant G$ is a nilpotent Hall $\pi$-subgroup, then there exists $x \in G$ such that $H \cap H^x = O_\pi(G)$.
\end{theorem}

\begin{lemma} \label{lem:key}
Let $G$ be a group of odd order where $F(G)$ is a Hall subgroup and $G/F(G)$ is nilpotent.
Let $(P_i)_{i=1}^r$ be non-normal Sylow subgroups of $G$ for distinct primes.
If $A \unlhd G$ is abelian and $AN_G(P_i)=G$ for all $i$,
then there exists $x \in A$ such that $P_i \cap P_i^x =1$ for all $i$.
\end{lemma}

\begin{proof}
Let $F=F(G)$, and observe that $A \subseteq F$.
By the Schur-Zassenhaus theorem, there exists $K<G$ such that $G= F \rtimes K$.
We may assume that $\pi(G/F) = \{p_1,\ldots,p_r\}$ and that $P_i \in \Syl_{p_i}(G)$ for each $i$.
For fixed $i$, the hypothesis $A N_G(P_i)=G$ means that $A$ is transitive on $\Syl_{p_i}(G)$ by conjugation.
We first use this fact to show that $F(AK)=A$.

Let $K=K_1 \times \cdots \times K_r$ where $K_i \in \Syl_{p_i}(K) \subseteq \Syl_{p_i}(G)$.
If $g \in F$, then
$$ (AK)^g = AK^g = A(K_1^g \times \cdots \times K_r^g) . $$
By the transitivity of $A$ there exist $a_1,\ldots,a_r \in A$ such that
$$ (AK)^g = A(K_1^{a_1} \times \cdots \times K_r^{a_r}) . $$
Therefore $K_i^{a_i} \subseteq (AK)^g$ for all $i$, and since $a_i \in A$,
also $K_i \subseteq (AK)^g$ for all $i$.
This implies that $K \subseteq (AK)^g$ and so $AK \unlhd G$ because $g$ is arbitrary.
It follows that $F(AK) \subseteq AK \cap F = A$ and equality holds.

By Theorem \ref{thBialoHall} there exists $x \in AK$ such that $K \cap K^x \subseteq K \cap A =1$.
We can assume $x\in A$.
Of course we have
$$ K_i \cap K_i^x \subseteq K \cap K^x = 1 $$
for each $i$.
Using again the transitivity of $A$ on $\Syl_{p_i}(G)$, for each $i$ there exists $b_i\in A$ such that $P_i=K_i^{b_i}$.
But $A$ is abelian and so $b_i x = x b_i$ which gives
$$ P_i \cap P_i^x = K_i^{b_i} \cap K_i^{b_i x} = (K_i \cap K_i^x)^{b_i} = 1 $$
for all $i$, as desired.
\end{proof}

\end{document}
